\scanpage{54}
(\unsure{16}) (18), $V_{\mathbb C}\simeq\bigoplus_{\omega\in\underline\omega}V_{\mathbb C}^{\omega}$, le groupe produit $\mu_m^{\underline\omega}$ d'ordre $m^2$ opère sur $V_{\mathbb C}$, par
\[
(32)\qquad(\lambda_\omega)_{\omega\in\underline\omega}\cdot(x_\omega)_{\omega\in\underline\omega}=(\lambda_\omega x_\omega)
\]
On \deleted{\unsure{va}} a des hom
\[
(33)\qquad\left\{\begin{aligned}
\mu_m&\longrightarrow\mu_m^{\underline\omega}\qquad\text{hom diagonal}\\
\mu_m(\omega)\simeq D^+&\longrightarrow\mu_m^{\underline\omega}\\
\lambda&\longmapsto(\unsure{\lambda^\omega|\lambda_\omega})_{\omega\in\underline\omega}\quad\deleted{\unsure{\mathrm{dont}}}
\end{aligned}\right.
\]
tels que les actions de $\mu_m,D^+$ sur $V^{\deleted{\unsure{\omega}}}$ \deleted{\unsure{passent}} passent via celle de $\mu_m^{\underline\omega}$. L'hom correspondant
\[
(34)\qquad\deleted{\unsure{\mu_m}}\underset{\mu_m(\underline\omega)}{\underbrace{D^+}}\times\mu_m\longrightarrow\mu_m^{\underline\omega}
\]
entre groupes de $=$ ordre $m^2=4n^2$ a comme noyau le sous-groupe d'ordre $2$ engendré par \deleted{$\unsure{a_0,-1}$} $(a_0,-1)$, son image est donc d'indice $2$, elle est formée des couples $(\lambda_\omega)_{\omega\in\underline\omega}$ tels que l'on ait
\[
(35)\qquad\deleted{\unsure{\prod\lambda_\omega}}\left(\prod_{\omega\in\underline\omega}\lambda_\omega\right)^n=1\qquad\mathrm{i.e.}\quad(\lambda_\omega\lambda_{\omega'})^n=1.
\]
Notons que les autom. de $V_{\mathbb C}$ qui stabilisent $V_{\mathbb C}^{**}$ i.e. qui induisent un autom de $\Sigma$ qui stabilise $J$, sont ceux qui provien-
\mspage{53}
